\section{Limitations}
\label{sec:limitations}

\paragraph{Moderate tier accuracy on expert labels.}
The system achieves 0.53 tier accuracy on expert-labeled data, meaning
roughly half of all predictions match the exact tier assigned by a human
expert. While adjacent accuracy (0.84) demonstrates that most errors
are near-misses on an ordinal scale, exact-match performance remains a
barrier to fully unsupervised deployment. The primary error mode is
over-prediction of the Related tier, which acts as an attractor for
both Equivalent and Partial pairs.

\paragraph{Partial-tier confusion.}
The Partial class achieves the lowest per-class F1 (0.364) and accuracy
(0.308). The boundary between Partial and its neighbors (Unrelated and
Related) is inherently subjective: reasonable experts disagree on
whether a control ``partially addresses'' another versus being merely
``related to'' it. This ambiguity is reflected in the confusion matrix,
where 35/78 Partial pairs are predicted as Related and 19/78 as
Unrelated. Improving Partial discrimination likely requires richer
annotation guidelines or multi-annotator consensus labels.

\paragraph{Expert--LLM label distribution gap.}
The 12-point accuracy gap between LLM-labeled validation (${\sim}0.65$)
and expert-labeled test (0.53) indicates that the LLM-SME labeling
protocol produces systematically different tier distributions from human
experts, particularly for the Partial and Equivalent tiers. Training on
LLM labels while evaluating on expert labels introduces a distribution
shift that limits classifier performance.

\paragraph{Conservative conformal sets.}
Mondrian conformal prediction achieves 93.8\% coverage, exceeding the
nominal 90\% target, but at the cost of large prediction sets (mean
size 3.20 out of 4 classes). This conservatism limits the practical
utility of conformal prediction for downstream decision-making: in most
cases, the prediction set does not meaningfully narrow the tier
possibilities. Tighter sets would require either a more discriminative
base classifier or a larger calibration set.

\paragraph{Sparse framework connectivity.}
Only 15 of the 36 possible framework pairs (41.7\%) have cross-framework
edges in the current graph. The remaining 21 pairs have no direct
mapping evidence, meaning the classifier must rely entirely on text
similarity and Opus judgment for those pairs. Expanding upstream data
sources or applying zero-shot graph construction methods would improve
coverage.

\paragraph{Framework coverage and generalization.}
The current system covers nine AI security frameworks. Generalization to
new frameworks requires ingestion of control texts and ideally some
anchor edges, though the Opus tier-score features provide a framework-agnostic
baseline signal. Evaluation on out-of-distribution frameworks has not
been conducted.

\paragraph{Temporal validity.}
Framework versions are pinned at ingestion time. As frameworks evolve,
the crosswalk will require periodic re-ingestion, re-featurization, and
re-evaluation. The modular pipeline design facilitates this, but the
cost of re-labeling expert test sets for updated frameworks is
non-trivial.
